\documentclass[a4paper,11pt]{article}

% ── Packages ──
\usepackage[utf8]{inputenc}
\usepackage[T1]{fontenc}
\usepackage{amsmath, amssymb, amsthm}
\usepackage{algorithm}
\usepackage{algpseudocode}
\usepackage{booktabs}
\usepackage{graphicx}
\usepackage{hyperref}
\usepackage[margin=2.5cm]{geometry}
\usepackage{xcolor}
\usepackage{listings}

\hypersetup{
    colorlinks=true,
    linkcolor=blue!70!black,
    citecolor=green!50!black,
    urlcolor=blue!60!black
}

\lstset{
    language=Python,
    basicstyle=\ttfamily\small,
    keywordstyle=\color{blue!70!black}\bfseries,
    commentstyle=\color{green!50!black},
    stringstyle=\color{red!60!black},
    frame=single,
    breaklines=true
}

\title{%
    \textbf{Algorithm Documentation: Parallel Base Optimizer} \\[6pt]
    \large Average Manipulability Index and Grid Search Strategy \\
    for TM5-700 Surgical Robot Base Optimization
}
\author{RMML Lab — Isaac Lab Reachability Map Project}
\date{\today}

\begin{document}
\maketitle

% ============================================================
\section{Overview}
% ============================================================

The \texttt{parallel\_base\_optimizer.py} script performs a \textbf{2D grid search} over candidate robot base positions using Isaac Lab's parallel environment framework. For each candidate position, it evaluates the robot's ability to reach a set of surgical waypoints and computes a kinematic quality metric — the \textbf{Average Manipulability Index (AMI)} — to determine the optimal base placement for a TM5-700 robotic arm in a surgical scenario.

% ============================================================
\section{Coordinate System and Base Positioning}
% ============================================================

\subsection{USD Hierarchy}

The robot exists within a nested USD (Universal Scene Description) hierarchy:

\begin{verbatim}
/Root
  |-- Zero                        (global origin reference)
  +-- robotarm_base               (grid search target)
       +-- robotarm_base          (second-level Xform)
            +-- tm5_700           (articulation root)
\end{verbatim}

Isaac Lab controls the \textbf{articulation root} (\texttt{tm5\_700}) via physics APIs, while the user specifies the desired position of \texttt{robotarm\_base} (first level). A fixed offset $\mathbf{T}$ relates the two:

\begin{equation}
    \mathbf{p}_{\text{tm5}} = \mathbf{p}_{\text{base}} + \mathbf{T}
    \label{eq:offset}
\end{equation}

where $\mathbf{T} \in \mathbb{R}^3$ is computed at runtime:

\begin{equation}
    \mathbf{T} = \mathbf{p}_{\text{tm5}}^{\text{default}} - \mathbf{p}_{\text{base}}^{\text{default}}
\end{equation}

For the current USD scene, $\mathbf{p}_{\text{base}}^{\text{default}} = (0.2, 0, 0)^T$ m.

\subsection{Grid Search Definition}

The search space is defined as a 2D Cartesian grid:

\begin{equation}
    \mathcal{G} = \{(x_i, y_j) \mid x_i \in \mathtt{linspace}(x_{\min}, x_{\max}, N_x),\ y_j \in \mathtt{linspace}(y_{\min}, y_{\max}, N_y)\}
\end{equation}

Each grid point $(x_i, y_j)$ represents a candidate \texttt{robotarm\_base} position. The total number of parallel environments is $|\mathcal{G}| = N_x \times N_y$.

% ============================================================
\section{Jacobian Matrix and Manipulability Index}
\label{sec:manipulability}
% ============================================================

\subsection{Jacobian Matrix}

For a 6-DOF serial manipulator with joint configuration $\mathbf{q} \in \mathbb{R}^6$, the \textbf{geometric Jacobian} $\mathbf{J}(\mathbf{q}) \in \mathbb{R}^{6 \times 6}$ relates joint velocities to end-effector (needle tip) velocities:

\begin{equation}
    \dot{\mathbf{x}} = \mathbf{J}(\mathbf{q})\, \dot{\mathbf{q}}
    \label{eq:jacobian}
\end{equation}

where $\dot{\mathbf{x}} = [\mathbf{v}^T, \boldsymbol{\omega}^T]^T$ contains the linear and angular velocities of the end-effector.

In the script, the Jacobian is extracted directly from the PhysX simulation:

\begin{lstlisting}
J = robot.root_physx_view.get_jacobians()
    [:, ee_jacobi_idx, :, joint_col_ids]
\end{lstlisting}

This retrieves the $6 \times 6$ sub-matrix corresponding to the needle tip body and the 6 active joints (\texttt{joint\_1} through \texttt{joint\_6}).

\subsection{Yoshikawa Manipulability Index}

The \textbf{Manipulability Index} $w$, introduced by Yoshikawa~(1985)\cite{yoshikawa1985}, is defined as:

\begin{equation}
    \boxed{w(\mathbf{q}) = \left| \det\bigl(\mathbf{J}(\mathbf{q})\bigr) \right|}
    \label{eq:manip}
\end{equation}

\subsubsection{Geometric Interpretation}

The value $w(\mathbf{q})$ represents the \textbf{volume of the velocity ellipsoid} — the set of all achievable end-effector velocities for unit joint velocities:

\begin{equation}
    \mathcal{E}(\mathbf{q}) = \left\{ \dot{\mathbf{x}} \in \mathbb{R}^6 \;\middle|\; \|\dot{\mathbf{q}}\| \leq 1,\; \dot{\mathbf{x}} = \mathbf{J}(\mathbf{q})\dot{\mathbf{q}} \right\}
\end{equation}

\begin{itemize}
    \item $w \to 0$: The ellipsoid collapses along one or more axes $\Rightarrow$ the robot is near a \textbf{kinematic singularity}, losing the ability to move in certain Cartesian directions.
    \item $w$ large: The ellipsoid is ``full'' $\Rightarrow$ the robot can move and rotate the end-effector effectively in all directions — high \textbf{dexterity}.
\end{itemize}

\subsubsection{Relation to Singular Values}

The manipulability index can also be expressed in terms of the singular values $\sigma_i$ of $\mathbf{J}$:

\begin{equation}
    w = \prod_{i=1}^{6} \sigma_i
\end{equation}

If any $\sigma_i = 0$, the robot is at a singularity and $w = 0$.

% ============================================================
\section{Average Manipulability Index (AMI)}
\label{sec:ami}
% ============================================================

For a given base position $\mathbf{p}_k = (x_k, y_k)$ and a set of $M$ task waypoints $\{\mathbf{w}_1, \ldots, \mathbf{w}_M\}$, define:

\begin{itemize}
    \item $\mathcal{S}_k \subseteq \{1, \ldots, M\}$: the set of waypoints \textbf{successfully reached} (positioning error $< \epsilon$, where $\epsilon = 1.5$ mm by default).
    \item $\mathbf{q}_k^{(j)}$: the joint configuration when waypoint $j \in \mathcal{S}_k$ is reached.
\end{itemize}

The \textbf{Average Manipulability Index} for base position $k$ is:

\begin{equation}
    \boxed{\text{AMI}_k = \frac{1}{|\mathcal{S}_k|} \sum_{j \in \mathcal{S}_k} w\!\left(\mathbf{q}_k^{(j)}\right) = \frac{1}{|\mathcal{S}_k|} \sum_{j \in \mathcal{S}_k} \left|\det\!\left(\mathbf{J}\!\left(\mathbf{q}_k^{(j)}\right)\right)\right|}
    \label{eq:ami}
\end{equation}

The script also records the \textbf{minimum manipulability}:

\begin{equation}
    w_{\min,k} = \min_{j \in \mathcal{S}_k} w\!\left(\mathbf{q}_k^{(j)}\right)
    \label{eq:wmin}
\end{equation}

This captures the worst-case dexterity across all reachable waypoints.

\subsection{Interpretation for Base Selection}

The optimal robot base position $\mathbf{p}^* \in \mathcal{G}$ is chosen to maximize a composite fitness score ($fg$), ensuring both strict reachability across all task waypoints and continuous high dexterity along the trajectory:

\begin{equation}
    \mathbf{p}^* = \arg\max_{\mathbf{p}_{i,j}} \left[ fg \right]
\end{equation}

Where the fitness score $fg$ for a given base position $\mathbf{p}_{i,j}$ is defined as:

\begin{equation}
    fg = S_r^{min} \cdot \left( \prod_{k=1}^{m} c(t, k) \right) \cdot cs
\end{equation}

\begin{itemize}
    \item $S_r^{min} = \min_{k} (RI_k)$ : The minimum \textbf{Reachability Index (RI)} across all $m$ target waypoints for the base position. $RI_k = \frac{\text{\# successful orientations}}{\text{total orientations attempted}}$ for waypoint $k$. This acts as a strict penalty constraint; if any single waypoint is completely unreachable ($RI_k = 0$), the entire base position is deemed invalid ($fg = 0$), ensuring no dead-zones exist in the surgical task space.
    
    \item $c(t, k) = \text{AMI}_k$ : The dexterity cost function for the $k$-th waypoint ($1 \le k \le m$), defined as the \textbf{Average Manipulability Index (AMI)} across all \textbf{successful} local orientations:
    \begin{equation}
        \text{AMI}_k = \frac{1}{|\mathcal{S}_k|} \sum_{\text{successful orientations}} w(\mathbf{q})
    \end{equation}
    where $\mathcal{S}_k$ is the set of successful joint configurations (i.e., those reaching waypoint $k$ within error tolerance $\epsilon$). This represents the average dexterity only among reachable configurations; unreachable configurations contribute zero.
    
    \item $cs = \frac{\text{\# successful tasks}}{\text{total tasks}}$ : A scaling coefficient (success rate across all waypoints and orientations) to balance the overall fitness.
\end{itemize}

\subsubsection{Terminology Clarification: RI vs MI vs AMI}

\begin{itemize}
    \item \textbf{Manipulability Index (MI, $w$)}: Defined per joint configuration $\mathbf{q}$ as $w(\mathbf{q}) = |\det(\mathbf{J}(\mathbf{q}))|$. Scalar value $\ge 0$.
    \item \textbf{Reachability Index (RI, $RI_k$)}: Defined per waypoint $k$ as the fraction of local orientations that successfully reach the waypoint. Range $[0, 1]$.
    \item \textbf{Average Manipulability Index (AMI, $\text{AMI}_k$)}: Defined per waypoint $k$ as the average MI across only the successful joint configurations for that waypoint. Represents dexterity quality among reachable poses.
\end{itemize}

The product $\prod_{k=1}^{m} c(t, k) = \prod_{k=1}^{m} \text{AMI}_k$ aggregates dexterity across the entire trajectory, penalizing waypoints with low average manipulability even if they are technically reachable (high $RI_k$).

In cases where multiple base positions yield identical $fg$ scores (e.g., highly symmetric configurations or capped metrics), the secondary tie-breaker objective is to maximize the Average Manipulability Index across the entire workspace ($Avg\_w$):

\begin{equation}
    \text{If } fg(\mathbf{p}_A) = fg(\mathbf{p}_B), \text{ select } \arg\max (Avg\_w)
\end{equation}

% ============================================================
\section{Algorithm}
% ============================================================

\begin{algorithm}[H]
\caption{Parallel Base Position Optimization}
\label{alg:optimizer}
\begin{algorithmic}[1]
\Require Grid parameters $(x_{\min}, x_{\max}, N_x, y_{\min}, y_{\max}, N_y)$, waypoints $\{\mathbf{w}_j\}_{j=1}^M$, threshold $\epsilon$
\Ensure $\text{AMI}_k$, $|\mathcal{S}_k|$ for each grid point $k$

\State $\mathcal{G} \leftarrow \text{meshgrid}(x_{\min}:x_{\max}, y_{\min}:y_{\max})$
\State Create $|\mathcal{G}|$ parallel environments in Isaac Lab
\State $\mathbf{T} \leftarrow \mathbf{p}_{\text{tm5}}^{\text{default}} - \mathbf{p}_{\text{base}}^{\text{default}}$ \Comment{Fixed offset}

\For{each env $k = 0, \ldots, |\mathcal{G}|-1$}
    \State Set $\mathbf{p}_{\text{tm5}}^{(k)} \leftarrow \mathbf{p}_{\text{env\_origin}}^{(k)} + \mathbf{p}_k + \mathbf{T}$ \Comment{Eq.~\ref{eq:offset}}
\EndFor

\State \Call{write\_root\_pose\_to\_sim}{} \Comment{Single GPU-compatible call}
\State Settle joints to initial configuration (50 steps)

\For{each waypoint $j = 1, \ldots, M$}
    \State Compute target $\mathbf{w}_j^{(k)}$ for each env (shifted by env origin)
    \For{step $= 1, \ldots, \text{max\_steps}$}
        \State $\mathbf{J}^{(k)} \leftarrow$ \Call{get\_jacobians}{} for all envs
        \State $\Delta\mathbf{q}^{(k)} \leftarrow$ \Call{DifferentialIK}{$\mathbf{w}_j^{(k)}$, $\mathbf{J}^{(k)}$}
        \State Apply joint position targets with clamping
        \State \Call{sim.step}{}
        \If{all envs converged ($\|\mathbf{e}\| < \epsilon$)}
            \State \textbf{break}
        \EndIf
    \EndFor
    \State $d^{(k)} \leftarrow \|\mathbf{p}_{\text{ee}}^{(k)} - \mathbf{w}_j^{(k)}\|$
    \If{$d^{(k)} \leq \epsilon$}
        \State $|\mathcal{S}_k| \mathrel{+}= 1$
        \State $\text{manip\_sum}_k \mathrel{+}= |\det(\mathbf{J}^{(k)})|$
    \EndIf
\EndFor

\State $\text{AMI}_k \leftarrow \text{manip\_sum}_k \;/\; |\mathcal{S}_k|$ \Comment{Eq.~\ref{eq:ami}}
\State Export CSV and generate heatmap
\end{algorithmic}
\end{algorithm}

% ============================================================
\section{Differential IK Controller}
% ============================================================

The script uses the \textbf{Damped Least Squares (DLS)} variant of differential inverse kinematics:

\begin{equation}
    \Delta\mathbf{q} = \mathbf{J}^T \left(\mathbf{J}\mathbf{J}^T + \lambda^2 \mathbf{I}\right)^{-1} \Delta\mathbf{x}
    \label{eq:dls}
\end{equation}

where $\lambda = 0.01$ is the damping factor and $\Delta\mathbf{x} = \mathbf{x}_{\text{target}} - \mathbf{x}_{\text{current}}$ is the Cartesian error. This avoids numerical instability near singularities compared to the pseudo-inverse method.

An integral correction term is added when the end-effector is close to the target ($\|\Delta\mathbf{x}\| < 30$ mm):

\begin{equation}
    \Delta\mathbf{x}_{\text{boosted}} = \Delta\mathbf{x} + \mathbf{I}_{\text{cart}}, \quad \mathbf{I}_{\text{cart}} = \text{clamp}\!\left(\sum_{t} 0.02 \cdot \Delta\mathbf{x}_t,\ -0.05,\ 0.05\right)
\end{equation}

This acts as a proportional-integral (PI) controller in Cartesian space, improving steady-state accuracy.

% ============================================================
\section{Joint Clamping}
% ============================================================

After computing $\Delta\mathbf{q}$, joint position targets are clamped to prevent unsafe configurations:

\begin{equation}
    q_i^{\text{target}} = \text{clamp}(q_i^{\text{target}},\ q_i^{\min},\ q_i^{\max})
\end{equation}

\begin{table}[h]
\centering
\caption{Joint limits applied in the optimizer}
\begin{tabular}{ccc}
\toprule
\textbf{Joint} & $q_{\min}$ (rad) & $q_{\max}$ (rad) \\
\midrule
joint\_1 & $-3.14$ & $-0.1$ \\
joint\_2 & $-0.3$  & $2.5$  \\
joint\_3 & $-0.3$  & $2.5$  \\
joint\_4 & $-0.5$  & $3.14$ \\
joint\_5 & $-1.57$ & $1.57$ \\
joint\_6 & $-3.14$ & $0.5$  \\
\bottomrule
\end{tabular}
\label{tab:joints}
\end{table}

These limits are task-specific constraints (not hardware limits) designed to keep the arm in a feasible workspace for the surgical scenario.

% ============================================================
\begin{thebibliography}{9}

\bibitem{yoshikawa1985}
T.~Yoshikawa, ``Manipulability of robotic mechanisms,''
\textit{The International Journal of Robotics Research}, vol.~4, no.~2, pp.~3--9, 1985.

\bibitem{sundaram2022}
C.~P.~Sundaram \textit{et al.}, ``Robot-Assisted Surgical System Configuration Mapping and Procedure Planning,''
referenced in project plan for RASSCMAP cost function and reachability analysis, 2022.

\end{thebibliography}

\end{document}
